% GENERATED FILE. Edit canonical lessons, then run npm run generate:books.
% canonical-source-hash: fnv1a64:ce39c098aba21f81
% canonical-lessons: IT-C02-prego, IT-C02-prego-go-ahead, IT-C02-prego-here-you-are, IT-C02-prego-come-again, IT-C02-come, IT-C02-stare, IT-C02-come-stai, IT-C02-come-sta, IT-C02-come-va, IT-C02-cosi-cosi, IT-C02-non-ce-male, IT-C02-practice

\chapter{How Are You?}
\label{ch:it-how-are-you}
\hlchaptermodality{2}

\begin{chapteropening}
\textbf{By the end of this chapter:} \emph{I can ask an Italian how they are, formally or informally, answer for myself, and hand the question back.}

The full Come stai? exchange run both ways — sta and Lei for respect, stai and tu for a friend — answered with sto bene or così così and closed with grazie / prego.
\end{chapteropening}

\section[prego]{prego — \textquotedblleft{}you're welcome,\textquotedblright{} literally \textquotedblleft{}I pray\textquotedblright{}}

\label{lesson:IT-C02-prego}



\begin{quote}
The \emph{grazie} lesson already promised you this word. Someone says \textbf{grazie}; you reply \textbf{prego}. Hold that one exchange first. Later lessons will add the word's other jobs one at a time.
\end{quote}

\begin{sounds}
\begin{itemize}
\raggedright
  \item \texttt{hard-g-before-o} — \textbf{prego} = \emph{PREH-go} with a \textbf{hard} g. (Italian g is soft only before \emph{e/i}; here the g is followed by \emph{o}, so it stays hard.)
\end{itemize}
\end{sounds}

\begin{cousinweb}
\textbf{prego} is \textquotedblleft{}I pray\textquotedblright{} — the \emph{io} (\textquotedblleft{}I\textquotedblright{}) form of \textbf{pregare}, \textquotedblleft{}to pray, to beg,\textquotedblright{} from Latin \textbf{precārī}. When you answer \emph{grazie} with \emph{prego}, you're softly saying \emph{\textquotedblleft{}I pray (it was nothing) / please, don't mention it.\textquotedblright{}}

That root \textbf{prec-} (\textquotedblleft{}pray, entreat\textquotedblright{}) gives English a surprising family:

\begin{itemize}
\raggedright
  \item \textbf{pray}, \textbf{prayer} — the direct line.
  \item \textbf{precarious} — literally \textquotedblleft{}held by \emph{prayer},\textquotedblright{} i.e. depending on someone's goodwill, therefore unstable.
  \item \textbf{deprecate} (\textquotedblleft{}pray \emph{away}, ward off\textquotedblright{}) and \textbf{imprecation} (\textquotedblleft{}a prayer \emph{against} — a curse\textquotedblright{}).
\end{itemize}

For now, keep one complete turn:

\begin{quote}
\textbf{Grazie. — Prego.} — \textquotedblleft{}Thank you. — You're welcome.\textquotedblright{}
\end{quote}
\end{cousinweb}

\subsection*{Your turn}
Say these aloud:

\begin{itemize}
\raggedright
  \item \textquotedblleft{}grazie\textquotedblright{} … \textquotedblleft{}prego\textquotedblright{} — the full exchange
  \item \textquotedblleft{}prego\textquotedblright{} then English \textquotedblleft{}pray, precarious\textquotedblright{} — one family
\end{itemize}

\begin{tcolorbox}[breakable,colback=teal!4,colframe=teal!35!black,title={Before you move on}]
What verb is \emph{prego} the \textquotedblleft{}I\textquotedblright{} form of? (\emph{pregare}, \textquotedblleft{}to pray.\textquotedblright{}) What English word means \textquotedblleft{}held by prayer\textquotedblright{} and so \textquotedblleft{}unstable\textquotedblright{}? (\emph{Precarious}.) What do you answer after \textbf{grazie}? (\textbf{Prego.}) Next, add one more use without losing this one.
\end{tcolorbox}
\section[prego — go ahead]{prego — \textquotedblleft{}please, go ahead\textquotedblright{}}

\label{lesson:IT-C02-prego-go-ahead}



\begin{quote}
Answer \textbf{grazie} first: \textbf{prego}. Keep that meaning. Now picture a doorway and add one nearby use.
\end{quote}

\subsection*{You'll want to know — One more job}
When you let someone pass or offer a seat, \textbf{prego} means \textbf{\textquotedblleft{}please, go ahead\textquotedblright{} / \textquotedblleft{}after you.\textquotedblright{}} The polite invitation still grows from \emph{pregare}, \textquotedblleft{}to pray or entreat.\textquotedblright{}

\begin{quote}
You hold the door. \textbf{Prego.} — \textquotedblleft{}After you.\textquotedblright{}
\end{quote}

Do not add another meaning yet. The same sound now owns two situations: after \textbf{grazie}, \textquotedblleft{}you're welcome\textquotedblright{}; at a doorway, \textquotedblleft{}go ahead.\textquotedblright{}

\subsection*{Your turn}
Say these aloud:

\begin{itemize}
\raggedright
  \item \textquotedblleft{}Grazie. — Prego.\textquotedblright{}
  \item \textquotedblleft{}Prego\textquotedblright{} while yielding a doorway
\end{itemize}

\begin{tcolorbox}[breakable,colback=teal!4,colframe=teal!35!black,title={Before you move on}]
At a doorway, what does \textbf{prego} mean? (\textbf{Please, go ahead / after you.}) After \textbf{grazie}? (\textbf{You're welcome.}) Keep both; the next lesson adds only one more.
\end{tcolorbox}
\section[prego — here you are]{prego — \textquotedblleft{}here you are\textquotedblright{}}

\label{lesson:IT-C02-prego-here-you-are}



\begin{quote}
At a doorway, \textbf{prego} invites a person forward. Move the same invitation into your hand.
\end{quote}

\subsection*{You'll want to know — One more job}
When you hand something to someone, \textbf{prego} means \textbf{\textquotedblleft{}here you are.\textquotedblright{}} The object can be a cup, a ticket, or anything already requested.

\begin{quote}
You pass over the cup. \textbf{Prego.} — \textquotedblleft{}Here you are.\textquotedblright{}
\end{quote}

The meanings stay tied to visible situations: thanks, doorway, offered object. You do not have to guess between them without context.

\subsection*{Your turn}
Say these aloud:

\begin{itemize}
\raggedright
  \item \textquotedblleft{}Prego\textquotedblright{} while yielding a doorway
  \item \textquotedblleft{}Prego\textquotedblright{} while handing over an imaginary cup
\end{itemize}

\begin{tcolorbox}[breakable,colback=teal!4,colframe=teal!35!black,title={Before you move on}]
You hand someone a cup. What does \textbf{prego} mean? (\textbf{Here you are.}) At a doorway? (\textbf{Go ahead.}) After thanks? (\textbf{You're welcome.})
\end{tcolorbox}
\section[Prego? — come again?]{Prego? — \textquotedblleft{}sorry, come again?\textquotedblright{}}

\label{lesson:IT-C02-prego-come-again}



\begin{quote}
You know \textbf{prego} after thanks, at a doorway, and with an offered object. The last common use changes the voice, not the spelling.
\end{quote}

\subsection*{You'll want to know — One more job}
If you did not hear, lift your voice into a question:

\begin{quote}
\textbf{Prego?} — \textquotedblleft{}Sorry? Come again?\textquotedblright{}
\end{quote}

The rising question sound separates this repair from a flat \textbf{prego} while handing something over. Listen to the situation and the voice together.

\subsection*{Your turn}
\begin{itemize}
\raggedright
  \item \emph{Say it:} flat \textquotedblleft{}Prego\textquotedblright{} — here you are
  \item \emph{Ask:} rising \textquotedblleft{}Prego?\textquotedblright{} — come again?
\end{itemize}

\begin{tcolorbox}[breakable,colback=teal!4,colframe=teal!35!black,title={Before you move on}]
You did not catch a sentence. What do you ask? (\textbf{Prego?}) What marks the listening-repair use? (\textbf{Question intonation and the situation.})
\end{tcolorbox}
\section[come]{come — \textquotedblleft{}how,\textquotedblright{} the third \emph{quomodo} cousin}

\label{lesson:IT-C02-come}



\begin{quote}
To ask \emph{how are you?} you need \textquotedblleft{}how.\textquotedblright{} In Italian it's \textbf{come} — and if you've seen the Spanish or French tracks, you already know its family.
\end{quote}

\begin{sounds}
\begin{itemize}
\raggedright
  \item \texttt{soft-c} — \textbf{come} = \emph{KOH-meh}. Here \emph{c} is before \emph{o}, so it's a \textbf{hard} \emph{k} (Italian \emph{c} only softens to \emph{ch} before \emph{e/i}). Final \emph{-e} is sounded: \emph{KOH-meh}.
\end{itemize}
\end{sounds}

\begin{cousinweb}
\textbf{come} = \textquotedblleft{}how / like / as.\textquotedblright{} Root: Latin \textbf{quōmodo}, \textquotedblleft{}in what manner\textquotedblright{} — two words (\emph{quō} \textquotedblleft{}in what way\textquotedblright{} + \emph{modo} \textquotedblleft{}manner\textquotedblright{}) fused and worn down. It is the \textbf{same source} as:

\begin{itemize}
\raggedright
  \item Spanish \textbf{cómo} (\textquotedblleft{}how\textquotedblright{})
  \item French \textbf{comment} (\textquotedblleft{}how\textquotedblright{})
  \item and, doubling as \textquotedblleft{}like/as,\textquotedblright{} it matches Spanish \textbf{como} and French \textbf{comme}.
\end{itemize}

So one Latin \emph{quōmodo} became the \textquotedblleft{}how\textquotedblright{} word in all three Romance sisters. The \emph{modo} piece (\textquotedblleft{}manner, measure\textquotedblright{}) also lives in English \textbf{mode}, \textbf{mod}el, \textbf{mod}erate — all about the \emph{way} something is done.

Note Italian \emph{come} covers \textbf{both} jobs at once, and the written accent tells them apart only by context, not spelling: \emph{Come stai?} (\textquotedblleft{}\textbf{How} are you?\textquotedblright{}) vs \emph{bianco come la neve} (\textquotedblleft{}white \textbf{as} snow\textquotedblright{}).
\end{cousinweb}

\subsection*{Your turn}
Say these aloud:

\begin{itemize}
\raggedright
  \item \textquotedblleft{}come\textquotedblright{} — \emph{KOH-meh}
  \item come / cómo / comment — the three \emph{quomodo} siblings
  \item English \emph{mode, model} share the \emph{modo} piece
\end{itemize}

\begin{tcolorbox}[breakable,colback=teal!4,colframe=teal!35!black,title={Before you move on}]
What Latin word gives \emph{come}, and which Spanish and French words are its siblings? (\emph{quōmodo} $\to$ \emph{cómo}, \emph{comment}.) What two jobs does \emph{come} do? (\textquotedblleft{}How?\textquotedblright{} and \textquotedblleft{}like/as.\textquotedblright{}) Next: the verb it pairs with — \textbf{stare}.
\end{tcolorbox}
\section[stare]{stare — \textquotedblleft{}to be,\textquotedblright{} the standing kind (Spanish's estar, undressed)}

\label{lesson:IT-C02-stare}



\begin{quote}
Like Spanish, Italian has \textbf{two} verbs for \textquotedblleft{}to be\textquotedblright{} — \textbf{essere} (identity) and \textbf{stare} (state). \textquotedblleft{}How are you?\textquotedblright{} runs on \textbf{stare}. Spanish's \emph{estar} is this very verb, with a front \emph{e} added on.
\end{quote}

\begin{sounds}
\begin{itemize}
\raggedright
  \item \texttt{st-clean} — Italian, unlike Spanish, is happy to start a word with \emph{st-}: no propping \emph{e}. \textbf{stare} = \emph{STAH-reh}, clean \emph{st}, final \emph{-e} sounded.
\end{itemize}
\end{sounds}

\begin{cousinweb}
\textbf{stare} comes straight from Latin \textbf{stāre}, \emph{\textquotedblleft{}to stand.\textquotedblright{}} Spanish grew the same verb into \textbf{estar} (with a propping \emph{e-}); Italian kept it bare as \textbf{stare}. Same Latin root, same job: the \emph{temporary state} be-verb — how you \emph{stand} right now.

That \textbf{stā-} root stands behind a wall of English words: \textbf{stay}, \textbf{stand}, \textbf{state}, \textbf{status}, \textbf{stable}, \textbf{station}, \textbf{stance}, \textbf{e-state}, \textbf{con-stant}, \textbf{circum-stance}. So \emph{Come stai?} is, at the root, \textbf{\textquotedblleft{}how do you stand?\textquotedblright{}} — the very same picture as Spanish \emph{¿Cómo estás?}.
\end{cousinweb}

\begin{grammarlens}[title={Grammar Lens: essere vs stare, and the two forms you need}]
\noindent
\begin{tabularx}{\linewidth}{@{}>{\raggedright\arraybackslash}X>{\raggedright\arraybackslash}X>{\raggedright\arraybackslash}X@{}}
\toprule
\textbf{verb} & \textbf{from} & \textbf{job} \\
\midrule
\textbf{essere} & Latin \emph{esse} & identity — \emph{sono Marco}, \textquotedblleft{}I am Marco\textquotedblright{} \\
\textbf{stare} & Latin \emph{stāre} & state / how you are — \emph{sto bene}, \textquotedblleft{}I'm well\textquotedblright{} \\
\bottomrule
\end{tabularx}

For this chapter, just two forms of \emph{stare}:

\begin{itemize}
\raggedright
  \item \textbf{stai} — \textquotedblleft{}you are\textquotedblright{} (informal, \emph{tu})
  \item \textbf{sta} — \textquotedblleft{}you are\textquotedblright{} (formal, \emph{Lei}) / \textquotedblleft{}he/she is\textquotedblright{}
\end{itemize}
\end{grammarlens}

\subsection*{Your turn}
Say these aloud:

\begin{itemize}
\raggedright
  \item \textquotedblleft{}stare\textquotedblright{} — \emph{STAH-reh}, clean \emph{st}
  \item \textquotedblleft{}stai\textquotedblright{} … \textquotedblleft{}sta\textquotedblright{} — the two forms you'll ask with
  \item Italian \emph{stare} = Spanish \emph{estar}, same Latin \emph{stāre} \textquotedblleft{}to stand\textquotedblright{}
\end{itemize}

\begin{tcolorbox}[breakable,colback=teal!4,colframe=teal!35!black,title={Before you move on}]
What Latin verb is \emph{stare}, and which Spanish verb is its twin? (\textquotedblleft{}To stand,\textquotedblright{} \emph{stāre} — Spanish \emph{estar}.) Which Italian be-verb is for how you feel — \emph{essere} or \emph{stare}? (\emph{stare}.) Next: ask it — \textbf{Come stai?}
\end{tcolorbox}
\section[Come stai?]{Come stai? — “how are you?”}

\label{lesson:IT-C02-come-stai}



\begin{quote}
Both pieces are ready: \textbf{come} (“how”) and \textbf{stare / stai} (“to be,” from \emph{stand}). Assemble Italian’s everyday question for a friend.
\end{quote}

\begin{cousinweb}
\begin{quote}
\textbf{Come stai?} = “How are you?” (informal) literally: \textbf{“How do you stand?”}
\end{quote}

\begin{itemize}
\raggedright
  \item \textbf{come} = “how”
  \item \textbf{stai} = “you are / stand,” the \emph{tu} form of \emph{stare}
\end{itemize}

The subject \textbf{tu} (“you”) is already carried by \textbf{stai}, so Italian normally leaves it unspoken. Say \textbf{Come stai?}, not \emph{Come tu stai?}, unless you need to stress “you.” This is the same pro-drop habit you will use in later verb lessons.

An everyday answer reuses the same verb:

\begin{quote}
\textbf{Come stai? — Sto bene, grazie.} “How are you? — I’m well, thanks.”
\end{quote}

\textbf{Sto} means “I am / stand.” The pair \textbf{stai / sto} lets the ending do the pronoun’s work: “you” in \emph{-ai}, “I” in \emph{-o}.
\end{cousinweb}

\subsection*{Your turn}
Say these aloud:

\begin{itemize}
\raggedright
  \item “Come stai?” — informal, for a friend
  \item “Sto bene, grazie.”
  \item the two-part exchange — “Come stai? — Sto bene.”
\end{itemize}

\emph{Twice through:} “Come stai? — Sto bene, grazie.”

\begin{tcolorbox}[breakable,colback=teal!4,colframe=teal!35!black,title={Before you move on}]
\emph{Come stai?} literally asks what? (“How do you \textbf{stand}?”) Which word means “how”? (\textbf{Come}.) Which form answers “I am”? (\textbf{Sto}.) Give the whole exchange: \textbf{Come stai? — Sto bene, grazie.} Next: formal speech and the “go” version.
\end{tcolorbox}
\section[Come sta?]{Come sta? — the formal question}

\label{lesson:IT-C02-come-sta}



\begin{quote}
You can ask a friend \textbf{Come stai?} Change one ending when you speak to a stranger or elder.
\end{quote}

\begin{grammarlens}[title={Grammar Lens: stai and sta}]
\begin{quote}
\textbf{Come sta?} = “How are you?” (formal)
\end{quote}

Italian uses \textbf{Lei} — literally “she,” often capitalized — as a polite “you.” Its verb therefore uses the third-person form \textbf{sta}, not informal \textbf{stai}:

\noindent
\begin{tabularx}{\linewidth}{@{}>{\raggedright\arraybackslash}X>{\raggedright\arraybackslash}X@{}}
\toprule
\textbf{listener} & \textbf{question} \\
\midrule
a friend (\emph{tu}) & \textbf{Come stai?} \\
a stranger or elder (\emph{Lei}) & \textbf{Come sta?} \\
\bottomrule
\end{tabularx}

The politeness pattern has cousins in Spanish \emph{usted} and German \emph{Sie}, but you only need one Italian contrast here: \textbf{stai} for a friend, \textbf{sta} for formal respect.

An answer can return the formal pronoun:

\begin{quote}
\textbf{Bene, grazie. E Lei?} — “Well, thank you. And you?”
\end{quote}
\end{grammarlens}

\subsection*{Your turn}
Say these aloud:

\begin{itemize}
\raggedright
  \item to a friend — “Come stai?”
  \item to a stranger — “Come sta?”
  \item “Bene, grazie. E Lei?”
\end{itemize}

\emph{Twice through:} “Stai — friend. Sta — formal.”

\begin{tcolorbox}[breakable,colback=teal!4,colframe=teal!35!black,title={Before you move on}]
Which ending marks a friend? (\textbf{-i: stai}.) Which form is formal? (\textbf{Sta}.) What polite pronoun can follow \emph{e}, “and”? (\textbf{Lei}.) Next: keep the meaning but switch the picture from “stand” to “go.”
\end{tcolorbox}
\section[Come va?]{Come va? — “how is it going?”}

\label{lesson:IT-C02-come-va}



\begin{quote}
\textbf{Come stai?} pictures wellbeing as standing. Italian also has the same “going” picture as English.

\begin{quote}
\textbf{Come va?} = “How’s it going?”
\end{quote}

\textbf{Va} comes from \textbf{andare}, “to go.” Unlike \emph{stai / sta}, this question does not force you to choose an informal or formal ending, so it is useful in either setting.
\end{quote}

\begin{grammarlens}[title={Grammar Lens: standing and going}]
\begin{itemize}
\raggedright
  \item \textbf{Come stai?} — “How do you \textbf{stand}?”
  \item \textbf{Come va?} — “How does it \textbf{go}?”
\end{itemize}

Spanish also uses the “stand” picture in \emph{¿cómo estás?}; French and German use a “go” picture. Italian keeps \textbf{both}, which makes it a bridge between them. The comparison is only a memory aid: in Italian, simply connect \textbf{stai} with \emph{stare} and \textbf{va} with \emph{andare}.

You will learn all six forms of \emph{andare} much later. For now, \textbf{va} is one ready-made piece inside this question.
\end{grammarlens}

\subsection*{Your turn}
Say these aloud:

\begin{itemize}
\raggedright
  \item the standing question — “Come stai?”
  \item the going question — “Come va?”
  \item “Come va? — Bene, grazie.”
\end{itemize}

\emph{Twice through:} “Come stai? — Come va?”

\begin{tcolorbox}[breakable,colback=teal!4,colframe=teal!35!black,title={Before you move on}]
Which question uses “stand”? (\textbf{Come stai?}) Which uses “go”? (\textbf{Come va?}) What infinitive owns \textbf{va}? (\textbf{Andare}.) Must \emph{Come va?} choose between \emph{tu} and \emph{Lei}? (\textbf{No}.) Next: answer with the original Italian “so-so.”
\end{tcolorbox}
\section[così così]{così così — the original \textquotedblleft{}so-so\textquotedblright{}}

\label{lesson:IT-C02-cosi-cosi}



\begin{quote}
English \textquotedblleft{}so-so\textquotedblright{} is a \textbf{loan translation of the Italian} \emph{così così}. Someone asks \emph{Come stai?}; you're middling; you wobble a hand and say \emph{così così} — \textquotedblleft{}thus and thus,\textquotedblright{} this way and that.
\end{quote}

\begin{sounds}
\begin{itemize}
\raggedright
  \item \texttt{soft-c} — \textbf{così} = \emph{koh-ZEE}: the \emph{c} is before \emph{o} (hard \emph{k}), and the single \emph{s} between vowels is \textbf{voiced}, a \emph{z} sound: \emph{koh-\textbf{ZEE}}. Stress the end (that's what the accent on \emph{ì} marks).
\end{itemize}
\end{sounds}

\begin{cousinweb}
\textbf{così} = \textquotedblleft{}thus, so, in this way,\textquotedblright{} from Latin \textbf{(ec)cum sīc}, \textquotedblleft{}(behold) so.\textquotedblright{} Doubled — \emph{così così} — it means \textquotedblleft{}\textbf{so} and \textbf{so},\textquotedblright{} neither up nor down: a verbal shrug.

The core piece \textbf{sīc} (\textquotedblleft{}thus, so\textquotedblright{}) is one you already write in English:

\begin{itemize}
\raggedright
  \item \textbf{sic} — the \emph{{[}sic{]}} you put after a quoted mistake means \textquotedblleft{}\textbf{thus} it was written.\textquotedblright{}
  \item It's also the Latin behind Spanish/French/Italian \textbf{sì / si / sì} (\textquotedblleft{}yes\textquotedblright{} — \textquotedblleft{}it is \emph{so}\textquotedblright{}), and the \emph{-si} inside French \textbf{ain-si} (\textquotedblleft{}thus\textquotedblright{}).
\end{itemize}

So \emph{così così} is literally \textbf{\textquotedblleft{}so, so\textquotedblright{}} — and English simply borrowed the shrug whole.
\end{cousinweb}

\begin{grammarlens}[title={Grammar Lens: your answer ladder for \emph{Come stai?}}]
\noindent
\begin{tabularx}{\linewidth}{@{}>{\raggedright\arraybackslash}X>{\raggedright\arraybackslash}X@{}}
\toprule
\textbf{answer} & \textbf{sense} \\
\midrule
\textbf{(Molto) bene, grazie} & (very) well, thanks \\
\textbf{Così così} & so-so \\
\bottomrule
\end{tabularx}

The cross-track shrug set, now five deep: Italian \textbf{così così} (\textquotedblleft{}so, so\textquotedblright{}), French \textbf{comme ci comme ça} (\textquotedblleft{}like this, like that\textquotedblright{}), Spanish \textbf{más o menos} (\textquotedblleft{}more or less\textquotedblright{}), German \textbf{es geht} (\textquotedblleft{}it goes\textquotedblright{}).
\end{grammarlens}

\subsection*{Your turn}
Say these aloud:

\begin{itemize}
\raggedright
  \item \textquotedblleft{}così così\textquotedblright{} — \emph{koh-ZEE koh-ZEE}, with a hand-wobble
  \item answer \textquotedblleft{}Come stai?\textquotedblright{} two ways — \textquotedblleft{}Bene\textquotedblright{} / \textquotedblleft{}Così così\textquotedblright{}
  \item \textquotedblleft{}sic\textquotedblright{} in English — \textquotedblleft{}thus\textquotedblright{} — the same \emph{sīc}
\end{itemize}

\begin{tcolorbox}[breakable,colback=teal!4,colframe=teal!35!black,title={Before you move on}]
What does \emph{così così} literally mean? (\textquotedblleft{}Thus, thus / so, so.\textquotedblright{}) What Latin word for \textquotedblleft{}so\textquotedblright{} is inside it, and where do you still write it in English? (\emph{sīc} $\to$ \emph{{[}sic{]}}.) English \textquotedblleft{}so-so\textquotedblright{} came from — which language? (Italian \emph{così così}.) Next: put the exchange together.
\end{tcolorbox}
\section[non c'è male]{non c'è male — \textquotedblleft{}not bad\textquotedblright{}}

\label{lesson:IT-C02-non-ce-male}



\begin{quote}
You can answer \textbf{Come stai?} with \textbf{bene} or \textbf{così così}. Add one more answer, on its own step.
\end{quote}

\subsection*{You'll want to know — The formula}
\begin{quote}
\textbf{Non c'è male.} — \textquotedblleft{}Not bad.\textquotedblright{}
\end{quote}

Learn the four sounds as one useful answer for now. The pieces behind it — negation and the \textbf{c'è} construction — belong to later lessons. You do not need their rules to answer honestly today.

\subsection*{Your turn}
Say these aloud:

\begin{itemize}
\raggedright
  \item \textquotedblleft{}Come stai? — Così così.\textquotedblright{}
  \item \textquotedblleft{}Come stai? — Non c'è male.\textquotedblright{}
\end{itemize}

\begin{tcolorbox}[breakable,colback=teal!4,colframe=teal!35!black,title={Before you move on}]
Give the new \textquotedblleft{}not bad\textquotedblright{} answer. (\textbf{Non c'è male.}) Must you explain each word yet? (\textbf{No. It is an unanalyzed pre-A1 formula for now.})
\end{tcolorbox}
\section[Practice]{Practice — \textquotedblleft{}Come stai?\textquotedblright{}, start to finish}

\label{lesson:IT-C02-practice}



\begin{quote}
No new words. This chapter's atoms — \emph{prego, come, stare, come stai?, così così} — assemble into a real exchange, formal and informal. It picks up from the Chapter 1 greetings.
\end{quote}

\subsection*{The exchange}
\textbf{Formal} (a stranger, an elder — \emph{Lei}):

\begin{quote}
— Buongiorno. \textbf{Come sta?} — \textbf{Bene, grazie.} E Lei? — \textbf{Così così.} — {[}after a small favour{]} — Grazie! — \textbf{Prego.}
\end{quote}

\textbf{Informal} (a friend — \emph{tu}):

\begin{quote}
— Ciao! \textbf{Come stai?} — \textbf{Sto bene, e tu?} — Non c'è male.
\end{quote}

The glue is \textbf{e Lei? / e tu?} — \textquotedblleft{}and you?\textquotedblright{} (\emph{e} = \textquotedblleft{}and,\textquotedblright{} Latin \emph{et}) — the Italian twin of Spanish \emph{¿y usted?} and French \emph{et toi?}.

\begin{grammarlens}[title={What you've built this chapter}]
\begin{itemize}
\raggedright
  \item \textbf{prego} — you're welcome ($\leftarrow$ \emph{pregare}, \textquotedblleft{}to pray\textquotedblright{}; English \emph{precarious}).
  \item \textbf{come} — how ($\leftarrow$ \emph{quōmodo}; sibling of \emph{cómo} / \emph{comment}).
  \item \textbf{stare} — \textquotedblleft{}to be\textquotedblright{} from \emph{stand} (= Spanish \emph{estar}); vs \emph{essere} (identity).
  \item \textbf{Come stai? / sta? / va?} — informal, formal, and the \textquotedblleft{}how's it going?\textquotedblright{} form (\emph{va} $\leftarrow$ \emph{andare}, \textquotedblleft{}to go\textquotedblright{}).
  \item \textbf{così così} — so-so (English \textquotedblleft{}so-so\textquotedblright{} borrowed it whole).
\end{itemize}
\end{grammarlens}

\subsection*{Your turn}
Say these aloud:

\begin{itemize}
\raggedright
  \item the full \emph{formal} exchange, both voices
  \item the same \emph{informally} — \emph{sta $\to$ stai}, \emph{Lei $\to$ tu}
  \item all three questions — Come stai? / Come sta? / Come va?
\end{itemize}

\emph{Twice through:} \textquotedblleft{}Come stai? — Sto bene, grazie, e tu?\textquotedblright{} until it's reflex.

\begin{tcolorbox}[breakable,colback=teal!4,colframe=teal!35!black,title={Before you move on}]
Which verb carries \emph{Come stai?}, and what does it literally ask? (\emph{stare}, \textquotedblleft{}to stand\textquotedblright{} — \textquotedblleft{}how do you stand?\textquotedblright{}) What does \emph{prego} answer, and from what root? (\emph{grazie}; $\leftarrow$ \emph{pregare}, \textquotedblleft{}to pray.\textquotedblright{}) Italian can ask \textquotedblleft{}how are you\textquotedblright{} two ways — which verbs? (\emph{stare}, stand; \emph{andare} $\to$ \emph{va}, go.) Next chapter: \textbf{introducing yourself} — mi chiamo, come ti chiami.
\end{tcolorbox}
